\documentclass{article}
\PassOptionsToPackage{compress}{natbib}
\usepackage{neurips_2025}
\usepackage{graphicx} % Required for inserting images

\title{SMM}
\author{Roussel Rahman}
\author{Jeff Shrager}
\date{September 2025}

% Optional math and formatting packages
\usepackage{amsmath, amssymb}
\usepackage{graphicx}
\usepackage{booktabs}
\usepackage{xcolor}

\newcommand{\rrsuggestion}[1]{\textcolor{blue}{#1}}

% RR: Setting Boxes vs Colored texts for links, citations etc.

\usepackage{hyperref}
%\AtBeginDocument{
\hypersetup{
    colorlinks=true,
    linkcolor=blue,
    citecolor=teal, % ForestGreen, MidnightBlue, Magenta
    filecolor=green,
    urlcolor=blue}

\title{A Small Math Model: Reimplementing Siegler and Shrager's Addition Model in an LLM-Inspired Architecture}

\author{
  Your Name \\
  Your Institution \\
  \texttt{your.email@example.com}
}

\begin{document}

\maketitle

\begin{abstract}
\citeauthor{siegler1984strategychoices}'s influential 1984 "distribution of associations" strategy choice model (S\&S84) introduced concepts into mathematical cognition that anticipated modern neural architectures. Their confidence-thresholded, associative strength recall mechanism, along with the use of backup strategies and ecologically-based training data, presaged key principles underlying today's successful Large Language Models (LLMs). Here we describe a ``Small Math Model'' (SMM) that modernizes the S\&S84 framework using contemporary neural architecture principles, and extends the original model to include a preceding learning-to-count stage -- a stage that was "baked in" to S\&S84. Like LLMs, the SMM processes symbolic sequences and produces probability distributions over responses. Because the model includes two operators (counting and addition) it also has to learn the distinguish these, and which arguments to attend to. \rrsuggestion{Trained this way, the SMM acquires/demonstrates several hallmark behavior predicated by the strategy choice theory.} Early results demonstrate plausible developmental strategy conflict resolution patterns, while the LLM-inspired design enables detailed analysis of emergent representations. We describe the model (which is available open source, of course) and our direction, which involves adding an embedding layer for inputs and outputs that should enable the model to learn emergent properties of numbers.
\end{abstract}

\section{Introduction}
The development of mathematical cognition has long been understood through the lens of strategy choice theory \citep{siegler1984strategychoices, siegler1987perils, siegler1991microgenetic, shrager1998scads}. Children use multiple strategies---counting, retrieval, decomposition---and select among them probabilistically based on confidence and associative strength. \citeauthor{siegler1984strategychoices}'s framework emphasized mechanisms such as backup strategies and a ``feeling of knowing,'' concepts that align with probabilistic inference, uncertainty estimation, and fallback mechanisms in modern AI.

Meanwhile, Large Language Models (LLMs) have revolutionized natural language processing by scaling these same principles. LLMs produce probabilistic outputs, estimate confidence via entropy, and dynamically select strategies from massive learned corpora. This parallel motivates a re-examination of early cognitive frameworks through the lens of modern architectures.

We present the Small Math Model (SMM), a shallow (three-layer) neural network designed to explore the counting-addition conflict central to early math development. The SMM allows us to revisit classical questions in a tractable contemporary architecture, and leverage contemporary analytic tools. \rrsuggestion{This way, the SMM provides an ideal testbed to trace LLM reasoning abilities as they emerge and compare with human reasoning, as steps towards an integrated understanding of intelligent reasoning in complex environments.}

\section{SSM<->LMM Analysis}

\section{Architecture and Learning Dynamics}

Our Small Math Model (SMM) is designed as a pedagogical analogue of a modern Large Math or Language Model (LMM/LLM). Inputs are tokenized into three slots---the first operand, operator, and second operand---each of which is mapped into a learned embedding space. This mirrors the embedding step of LLMs, albeit without positional encoding. The concatenated embeddings form the input to the network, analogous to a short sequence of tokens in a transformer.

To approximate attention, the SMM includes a learned, per-dimension gating mechanism. Each input dimension is weighted by a sigmoid gate parameterized by trainable matrices; this allows the model to emphasize or suppress parts of the concatenated input. While not a full self-attention mechanism, this ``quasi-attention'' plays a similar role in controlling the flow of information. The gated inputs are then processed by a hidden layer, projected into an embedding space, and scored against learned output embeddings before a softmax over the vocabulary of possible answers. Training is performed with a cross-entropy objective, exactly as in autoregressive LLMs, though over a much smaller vocabulary (numbers 1--12).

A distinctive feature of the SMM is its incorporation of explicit \emph{finger counting} as an auxiliary reasoning strategy. Confidence is estimated from the entropy of the output distribution; when confidence falls below a threshold, the model defers to an external finger-counting procedure that generates the correct result step by step. The output of this routine is then fed back into training, analogous to the way LLMs can be augmented with external tools or chains-of-thought and subsequently distilled back into the base model.

Training follows a curriculum schedule in which problem difficulty and task type (counting vs. addition) are varied according to Gaussian flows. This is directly inspired by curriculum learning practices in LLM training, where data distributions are staged over time to improve stability and generalization. Learning rates and confidence thresholds are annealed across steps, further paralleling hyperparameter scheduling in large-scale models.

In sum, the SMM modernizes and extends earlier work by Shrager et al.\ by explicitly drawing analogies to contemporary LLM architectures: embeddings, gating/attention, cross-entropy loss, curriculum learning, and tool use. While vastly simplified, this architecture provides a testbed for studying how counting priors, confidence estimation, and quasi-attention interact in a controlled setting, with the ultimate goal of bridging symbolic and neural perspectives on mathematical reasoning.


\section{Methods}
The SMM processes input sentences of the form ``a operator b'' where $a,b \in \{0,\ldots,9\}$ and operator $\in \{+\}$. Inputs are encoded as one-hot vectors and passed through a three-layer feedforward network trained with cross-entropy loss to predict the sum distribution.

\paragraph{Curriculum Learning.} Training follows a Gaussian curriculum schedule that gradually increases task difficulty. Early in training, simpler problems (e.g., sums with smaller operands) dominate; over time, harder problems appear with greater probability. This smooth progression mirrors natural learning trajectories while avoiding abrupt jumps.

\paragraph{Logging and Evaluation.} The implementation includes structured logging of accuracy, entropy, and error distributions across problem types. These diagnostics allow comparison to known developmental error patterns, such as over-reliance on counting or failure to generalize across contexts.

\section{Results}
Preliminary experiments demonstrate that the SMM reproduces several hallmark behaviors predicted by strategy choice theory:
\begin{itemize}
  \item \textbf{Counting-Addition Conflict:} The model initially overgeneralizes successor relationships (e.g., learning ``2$\to$3'' before ``2+3=5''), mirroring children's early developmental errors.
  \item \textbf{Entropy as Confidence:} Output entropy decreases systematically with training, reflecting increased certainty in well-practiced sums.
  \item \textbf{Backup Dynamics:} When addition accuracy falters, the model exhibits fallback to simpler successor-like behavior, analogous to procedural backup strategies.
\end{itemize}

Representative plots of accuracy vs. curriculum stage, entropy distributions, and confusion matrices highlight these emergent dynamics (to be included from simulation runs).

\rrsuggestion{Possible plots: Distributions of associative strength over all possible answers (similar to Fig. 9.8 from \citet{siegler1984strategychoices})}

\section{Discussion}
The SMM illustrates how modern neural implementations can instantiate and extend classical cognitive theories. By reframing Siegler and Shrager's ideas in computational terms aligned with LLMs, we obtain:
\begin{itemize}
  \item A platform for mechanistic analysis of symbolic reasoning in tractable domains.
  \item Opportunities to probe representations directly, e.g., via similarity analysis or attention visualization.
  \item A bridge between cognitive science and AI, highlighting shared principles of probabilistic selection, confidence estimation, and adaptive strategy use.
\end{itemize}

\section{Next Steps}
While the current implementation uses one-hot encodings, we plan to introduce \textbf{learned embeddings} for numbers and operators. Embeddings allow the model to discover emergent number properties (e.g., parity, proximity to 5 or 10) and support strategy-specific generalizations. To test this, we will design pseudo-strategies with error profiles dependent on number properties---for instance, a ``finger counting'' strategy that excels at $+1$ problems, or a $+5$ strategy that works best when an operand is 5. By varying error proneness we can examine whether embeddings capture such latent features. Probes (e.g., linear classifiers on embeddings) and behavioral splits (e.g., 1+ problems vs. 5+ problems) will serve as evaluation tools. This extension will provide a richer testbed for exploring how modern neural systems can internalize structured knowledge and reconcile competing strategies, further uniting insights from cognitive development and large-scale AI.

\section*{\rrsuggestion{Inventory -- TBD}}

\subsection*{Some high-level points}

(RR: Some points I thought important in general. Some may be too much to include within four pages.)

Investigating reasoning through aggregate scores obscure the use of different strategies, but true/accurate explanations need to specify the exact strategies \citep{newell1958elements, simon1971human, simon1976substantive, siegler1987perils, shrager1998scads, gigerenzer2020explain}. Reasoning and problem-solving in complex environments require complex skill organization, and it may be too complex to identify the elements of reasoning. A promising approach: explain reasoning in complex environments through the changes in strategies (i.e., their use) as new strategies are learned, in the same vein as \citet{siegler1984strategychoices} and \citet{shrager1998scads}.

The SMM forms the foundation of such an experimental testbed. A NN trained to learn addition by first learning to count and then building upon this ability to learn to add.

Catastrophic Forgetting vs Grokking. (What is grokking? \citet{nanda2023progress}) Catastrophic forgetting may stem from strict adherence to one strategy in uncertain situations that demand flexible use of different strategies. For example, in solving addition problems from memory, humans may revert to counting when confidence is low in the answer recalled from memory.

(what we may show here) The changes in strategies captured through the changes in step distributions. For humans, changes in strategy usage are reflected in changes in RT distributions. Here, we emulate the changes in the SMM. (specifics in a sentence) We show that .... shows the promise of the "microgenesis" of mathematical reasoning in environments of controllable complexity as a path to explaining AI reasoning in complex environments, such as mathematical problem-solving.

\subsection*{Some relevant research}
Catastrophic forgetting a well-established phenomenon in NNs, first observed by \citet{mccloskey1989catastrophic}
``New learning may interfere catastrophically with old learning when networks are trained sequentially'' \citep[][pp. 109]{mccloskey1989catastrophic}.

\bibliographystyle{apalike}
\bibliography{math_reasoning_bibliography}

\end{document}
